\subsection{Relatively invariant measures on a homogeneous space}

Let G again be a locally compact group, H a closed subgroup, \(\beta\) a left Haar measure on H.

Lemma 6. --- Let \(\lambda\) be a measure on G/H, \(\chi\) a continuous representation of G in C*. The following properties are equivalent:

\begin{enumerate}
\def\labelenumi{\alph{enumi})}
\item
  \(\lambda\) is relatively invariant on G/H with multiplier \(\chi\) ;
\item
  \(\lambda^{\sharp}\) is relatively invariant on G with left multiplier \(\chi\) ;
\item
  \(\lambda^{\sharp}\) is of the form \(a\chi \cdot \mu\) ( \(a \in \mathbf{C}\) ).
\end{enumerate}

The condition a) means that, for every \(s \in G\) ,

\[
\pmb {\gamma} _ {\mathrm{G/H}} (s) \lambda = \chi (s) ^ {- 1} \lambda ;
\]

this is equivalent to \(\left(\pmb{\gamma}_{\mathrm{G / H}}(s)\lambda\right)^{\sharp} = \chi (s)^{-1}\lambda^{\sharp}\) , that is, to

\[
\gamma_ {G} (s) \lambda^ {\sharp} = \chi (s) ^ {- 1} \lambda^ {\sharp}.
\]

Whence the equivalence of a) and b). The equivalence of b) and c) follows from §1, No.~8, Cor. 1 of Prop. 10.

THEOREM 3. --- Let G be a locally compact group, H a closed subgroup of G, \(\mu\) (resp. \(\beta\) ) a left Haar measure on G (resp. H), \(\chi\) a continuous representation of G in C*.

\begin{enumerate}
\def\labelenumi{\alph{enumi})}
\item
  In order that there exist on G/H a nonzero measure relatively invariant under G and with multiplier \(\chi\) , it is necessary and sufficient that \(\chi(\xi)=\Delta_{\mathrm{H}}(\xi)/\Delta_{\mathrm{G}}(\xi)\) for all \(\xi\in H\) .
\item
  This measure is then unique up to a constant factor; more precisely, it is proportional to \((\chi \cdot \mu) / \beta\) .
\end{enumerate}

For there to exist on G/H a nonzero measure relatively invariant under G with multiplier \(\chi\) , it is necessary and sufficient (Lemma 6) that \(\chi\cdot\mu\) be of the form \(\lambda^{\sharp}\) , hence (No.~2, Prop. 4) that \(\boldsymbol{\delta}(\xi)(\chi\cdot\mu)=\Delta_{\mathrm{H}}(\xi)(\chi\cdot\mu)\) for all \(\xi\in H\) . This condition may also be written \(\chi(\xi)\chi\cdot\Delta_{\mathrm{G}}(\xi)\mu=\Delta_{\mathrm{H}}(\xi)\chi\cdot\mu\) , that is,

\[
\chi (\xi) = \Delta_ {\mathrm{H}} (\xi) / \Delta_ {\mathrm{G}} (\xi)
\]

for all \(\xi \in \mathbf{H}\) . Whence a). The assertion b) follows at once from Lemma 6 and the fact that the mapping \(\lambda \mapsto \lambda^{\sharp}\) is injective.

\pandocbounded{\includegraphics[keepaspectratio]{images/part001_024ad5fd.jpg}}

We shall see in §3 (No.~3, Example 4) some very simple examples where the representation \(\xi \mapsto \Delta_{\mathrm{H}}(\xi) / \Delta_{\mathrm{G}}(\xi)\) cannot be extended to a continuous representation of G in \(\mathbf{C}^*\) . In this case, there therefore does not exist any nonzero complex measure on G/H relatively invariant under G.

COROLLARY 1. --- For there to exist on G/H a nonzero positive measure relatively invariant under G, it is necessary and sufficient that there exist a continuous representation of \(\mathbf{G}\) in \(\mathbf{R}_+^*\) extending the representation \(\xi \mapsto \Delta_{\mathrm{H}}(\xi) / \Delta_{\mathrm{G}}(\xi)\) .

Note that this condition is fulfilled when H is unimodular.

COROLLARY 2. --- For there to exist on G/H a nonzero positive measure invariant under G, it is necessary and sufficient that \(\Delta_{G}\) coincide with \(\Delta_{H}\) on H.

COROLLARY 3. --- Suppose that H is unimodular and that there exists on G/H a nonzero bounded positive measure \(\nu\) relatively invariant under G. Then \(\nu\) is invariant, and G is unimodular.

Let \(\chi\) be the multiplier of \(\nu\) . For every \(s \in G\) , \(\nu\) and \(\pmb{\gamma}(s)\nu\) have the same finite total mass (§1, No.~1, formula (6)); since \(\pmb{\gamma}(s)\nu = \chi(s)^{-1}\nu\) , we have \(\chi(s) = 1\) . Thus \(\nu\) is invariant. By Cor. 2, \(\Delta_G(s) = 1\) for all \(s \in H\) . Let \(G'\) be the set of \(t \in G\) such that \(\Delta_G(t) = 1\) . This is a closed normal subgroup of \(G\) containing \(H\) . Let \(\pi\) be the canonical mapping of \(G/H\) onto \(G/G'\) . Then \(\pi(\nu)\) is a nonzero, bounded positive measure invariant under \(G\) . Therefore the left Haar measure of the group \(G/G'\) is bounded, so that \(G/G'\) is compact (§1, No.~2, Prop. 2). Consequently the image of \(G\) under \(\Delta_G\) is a compact subgroup of \(\mathbf{R}_+^*\) ; this subgroup is reduced to \(\{1\}\) , thus \(\Delta_G = 1\) on all of \(G\) .

